% Part: computability
% Chapter: recursive-functions
% Section: partial-functions

\documentclass[../../../include/open-logic-section]{subfiles}

\begin{document}

\olfileid{cmp}{rec}{par}
\olsection{আংশিক পুনরাবৃত্ত অপেক্ষক}


পুনরাবৃত্ত অপেক্ষকের সংজ্ঞাটির প্রেরণা বুঝতে লক্ষ করুন, আদিম পুনরাবৃত্ত
নয় এমন গণনসাধ্য অপেক্ষক আছে—আমাদের প্রমাণটি আসলে আরও অনেক বেশি কিছু
প্রতিষ্ঠা করে। যুক্তিটি সরল ছিল: আমরা শুধু এই তথ্যটি ব্যবহার করেছিলাম যে
অপেক্ষকগুলিকে $f_0,f_1,\dots$ রূপে এমনভাবে তালিকায়িত করা যায়, যাতে
$x$ ও $y$-এর অপেক্ষক হিসেবে $f_x(y)$ গণনসাধ্য হয়। অতএব যুক্তিটি
\emph{এভাবে তালিকায়িত করা যায় এমন যেকোনো অপেক্ষক-শ্রেণি}-র ক্ষেত্রেই
প্রযোজ্য। এতে আমরা এক সংকটে পড়ি: গণনসাধ্য অপেক্ষকগুলিকে প্রত্যক্ষভাবে
বর্ণনা করতে চাই; কিন্তু গণনসাধ্য অপেক্ষকের কোনো সংগ্রহের প্রত্যক্ষ বর্ণনা
কখনোই সম্পূর্ণ হতে পারে না!

এই সংকট থেকে বেরোনোর উপায় হলো \emph{আংশিক} অপেক্ষককে আলোচনায় আনা।
আমরা দেখব যে আংশিক গণনসাধ্য অপেক্ষকগুলিকে তালিকায়িত করা সত্যিই
\emph{সম্ভব}।\iftag{TMs}{ আসলে, এটি যে সম্ভব তা আমরা প্রায় জেনেই গেছি,
  কারণ টুরিং যন্ত্রগুলিকে নিয়মমাফিক তালিকায়িত করা যায়।}{} পরে আমরা
আমাদের কর্ণ যুক্তিতে ফিরে দেখব, আংশিক অপেক্ষক অন্তর্ভুক্ত করলে যুক্তিটি
কেন আর কার্যকর হয় না।

এখন প্রশ্নটি হলো: সব আংশিক পুনরাবৃত্ত অপেক্ষক পেতে আদিম পুনরাবৃত্ত
অপেক্ষকগুলির সঙ্গে কী যোগ করতে হবে? আমাদের দুটি কাজ করতে হবে:
\begin{enumerate}
\item আদিম পুনরাবৃত্ত অপেক্ষকের সংজ্ঞাটিকে বদলে আংশিক অপেক্ষকেরও অনুমতি
  দিতে হবে।
\item সংজ্ঞাটিতে নতুন কিছু \emph{যোগ} করতে হবে, যাতে কিছু নতুন আংশিক
  অপেক্ষক অন্তর্ভুক্ত হয়।
\end{enumerate}

প্রথম কাজটি সহজ। আগের মতোই শূন্য, উত্তরসূরি ও অভিক্ষেপ দিয়ে শুরু করে
মিশ্রণ ও আদিম পুনরাবৃত্তির অধীনে বদ্ধতা নেব। একমাত্র পার্থক্য হলো,
সংজ্ঞার কয়েকটি পদ অসংজ্ঞায়িত হতে পারে—এই সম্ভাবনার অনুমতি দিতে মিশ্রণ
ও আদিম পুনরাবৃত্তির সংজ্ঞা বদলাতে হবে। $f$ ও $g$ আংশিক অপেক্ষক হলে,
$f(x) \fdefined$ দিয়ে বোঝাব যে $x$-এ $f$ সংজ্ঞায়িত, অর্থাৎ $x$ হলো
$f$-এর সংজ্ঞাক্ষেত্রের অন্তর্ভুক্ত; আর $f(x) \fundefined$ দিয়ে এর
বিপরীতটি, অর্থাৎ~$x$-এ $f$ অসংজ্ঞায়িত, বোঝাব। $f(x) \simeq g(x)$ দিয়ে
বোঝাব যে হয় $f(x)$ ও $g(x)$ উভয়ই অসংজ্ঞায়িত, নয়তো উভয়ই সংজ্ঞায়িত
ও সমান। আরও জটিল পদের ক্ষেত্রেও এই সংকেতলিপি ব্যবহার করব। আমরা এই
প্রথা নেব যে $h$ ও $g_0$, \dots,~$g_k$ সবাই আংশিক অপেক্ষক হলে
\[
h(g_0(\vec x),\dots,g_k(\vec x))
\]
সংজ্ঞায়িত হয় যদি এবং কেবল যদি প্রত্যেক $g_i$, $\vec x$-এ সংজ্ঞায়িত
হয় এবং $h$, $g_0(\vec x)$, \dots,~$g_k(\vec x)$-এ সংজ্ঞায়িত হয়। এই
বোঝাপড়া অনুযায়ী আংশিক অপেক্ষকের জন্য মিশ্রণ ও আদিম পুনরাবৃত্তির সংজ্ঞা
আগের মতোই, শুধু “$=$”-এর জায়গায় “$\simeq$” বসাতে হবে।

আংশিক অপেক্ষক পেতে আদিম পুনরাবৃত্ত অপেক্ষকের সংজ্ঞায় আমরা
\emph{সীমাহীন অনুসন্ধান অপারেটর} যোগ করব। স্বাভাবিক সংখ্যার উপর
$f(x,\vec z)$ যেকোনো আংশিক অপেক্ষক হলে, $\mu x \; f(x,\vec z)$-কে
নিচের বস্তু হিসেবে সংজ্ঞায়িত করি:
\begin{quote}
  এমন ক্ষুদ্রতম $x$, যার জন্য $f(0,\vec z), f(1,\vec z), \dots, f(x,\vec
  z)$ সবগুলি সংজ্ঞায়িত এবং $f(x,\vec z) = 0$—যদি এমন $x$ থাকে।
\end{quote}
অন্যথায় $\mu x \; f(x,\vec z)$ অসংজ্ঞায়িত হবে বলে বোঝা হবে। এতে
$\mu x \; f(x,\vec z)$ এককভাবে সংজ্ঞায়িত হয়।

\begin{explain}
লক্ষ করুন, আমাদের সংজ্ঞায়\iftag{TMs}{ টুরিং যন্ত্র, অথবা}{} অ্যালগরিদম
বা কোনো নির্দিষ্ট গণনামূলক মডেলের কোনো উল্লেখ নেই। কিন্তু মিশ্রণ ও আদিম
পুনরাবৃত্তির মতোই সীমাহীন অনুসন্ধানের পিছনে একটি প্রক্রিয়াগত, গণনামূলক
স্বজ্ঞা আছে। কোনো আংশিক অপেক্ষকের গণনসাধ্যতার ক্ষেত্রে, যেসব আর্গুমেন্টে
অপেক্ষকটি অসংজ্ঞায়িত সেগুলি এমন নিবেশের সঙ্গে সংশ্লিষ্ট, যেখানে গণনা
থামে না। $\mu x \; f(x,\vec z)$ গণনার পদ্ধতিটি হবে এই: শূন্য মান ফিরে না
আসা পর্যন্ত $f(0,\vec z), f(1,\vec z), f(2,\vec z)$ গণনা করতে থাকুন।
কিন্তু মধ্যবর্তী কোনো গণনা না থামলে $\mu x \; f(x,\vec z)$-এর গণনাও
থামবে না।
\end{explain}

$R(x,\vec z)$ যেকোনো সম্বন্ধ হলে, $\mu x \; R(x,\vec z)$-কে
$\mu x \; (1 \tsub \Char{R}(x,\vec z))$ হিসেবে সংজ্ঞায়িত করা হয়।
অন্যভাবে বললে, $\mu x \; R(x,\vec z)$ এমন ক্ষুদ্রতম $x$ ফেরত দেয় যার
জন্য $R(x,\vec z)$ সত্য। তাই $f(x,\vec z)$ সর্বত্র-সংজ্ঞায়িত অপেক্ষক
হলে $\mu x \; f(x,\vec z)$ এবং $\mu x \; (f(x,\vec z) = 0)$ একই।
কিন্তু লক্ষ করুন, আমাদের মূল সংজ্ঞাটি আরও সাধারণ, কারণ এতে $f(x,\vec z)$
সর্বত্র সংজ্ঞায়িত না-ও হতে পারে (অন্যদিকে কোনো সম্বন্ধের চরিত্রসূচক
অপেক্ষক সবসময় সর্বত্র-সংজ্ঞায়িত)।

\begin{defn}
স্বাভাবিক সংখ্যা থেকে স্বাভাবিক সংখ্যায় বিভিন্ন স্থানসংখ্যার আংশিক
অপেক্ষকগুলির যে ক্ষুদ্রতম সেট শূন্য, উত্তরসূরি ও অভিক্ষেপকে ধারণ করে এবং
মিশ্রণ, আদিম পুনরাবৃত্তি ও সীমাহীন অনুসন্ধানের অধীনে বদ্ধ, তাকে
\emph{আংশিক পুনরাবৃত্ত অপেক্ষকসমূহের} সেট বলা হয়।
\end{defn}

অবশ্য কিছু আংশিক পুনরাবৃত্ত অপেক্ষক সর্বত্র-সংজ্ঞায়িত হবে, অর্থাৎ
প্রত্যেক আর্গুমেন্টের জন্য সংজ্ঞায়িত হবে।

\begin{defn}
\ollabel{defn:recursive-fn}
\emph{পুনরাবৃত্ত অপেক্ষকসমূহের} সেট হলো সর্বত্র-সংজ্ঞায়িত আংশিক
পুনরাবৃত্ত অপেক্ষকগুলির সেট।
\end{defn}

কোনো পুনরাবৃত্ত অপেক্ষক যে সর্বত্র সংজ্ঞায়িত, তা জোর দিয়ে বোঝাতে একে
কখনও কখনও “সর্বত্র-সংজ্ঞায়িত পুনরাবৃত্ত” বলা হয়।

\end{document}
